% !TeX root = ../build/main.tex

% \noi Merkle trees are used to create a cryptographic representation of both the voter census registry and the election state. 

The state Merkle tree, denoted by $\hlset{\stateTree}$, represents the evolving global state of an election in \davinci. It is implemented as a \ac{SMT} of fixed depth $D = 64$, as described in Section~\ref{sec:cryptographic-primitives:merkle-trees}. The state tree acts as a cryptographic accumulator that compactly commits to all relevant election data, including configuration parameters, vote identifiers, and encrypted ballots.
By organizing the election state in a single \ac{SMT}, \davinci supports efficient membership and update proofs that can be verified inside \ac{ZK} circuits, enabling succinct and verifiable state transitions.

The root of the state tree, denoted $\stateroot$, is committed on-chain and serves as the authoritative reference to the current election state. Each state transition performed by a sequencer updates the tree and produces a new root. The correctness of this update, namely, that it follows the protocol rules and correctly incorporates new votes, is attested by a \ac{ZK-SNARK} proof, which is verified by the $\hlget{\statetransitionverifier}$ smart contract before the new root is accepted.


\paragraph{Namespaced key space.} 

The state tree operates over keys in the range $[0, 2^{64}-1] \subset \mathbb{F}_p$ and uses a SNARK-friendly hash function (Poseidon, see Section~\ref{sec:cryptographic-primitives:hash}) to enable efficient in-circuit verification. To prevent collisions between different categories of data stored in the state, the 64-bit key space is partitioned into three disjoint numeric namespaces using a fixed thershold $N$. We define
\[ N = 3 \cdot 2^{D-2},\]
which splits the address space into a lower region for configuration parameters, followed by a region for vote identifiers, and an upper region for encrypted ballot storage. This strict separation ensures that configuration entries, vote identifiers, and ballots can never collide. The structure of the tree is depicted in Figure~\ref{fig:mt-state}.

\begin{figure}[h]
	\centerline{\includegraphics[width=450pt,draft=false]{\figs/mt-state-mermaid}}
	\caption{Structure of the state Merkle tree. The first leaves are reserved for global election parameters while subsequent leaves store vote identifiers and encrypted ballots.}
	\label{fig:mt-state}
\end{figure}

\paragraph{Configuration namespace.}

The lowest indices of the tree are reserved for fixed process parameters that define the election rules. These keys are deterministic and occupy a negligible portion of the address space. Typical entries include:
\begin{itemize}
	\item \texttt{0x0}: the election identifier (\processid).
	\item \texttt{0x1}: the census origin (\censusorigin). %and \type, a field in case the organization was using a different structure for the census than a MT.
	\item \texttt{0x2}: the ballot mode configuration, which is the set of rules for validating votes (\ballotmode). 
	\item \texttt{0x3}: the encryption public key used to encrypt the ballots (\epk).
	\item \texttt{0x4}: an accumulator of the encrypted votes that need to be added to the tally (\addaccumulator). 
	\item \texttt{0x5}: an accumulator of the encrypted votes that have been overwritten (\substractaccumulator).
\end{itemize}
Let $\configmax$ denote the largest index reserved for configuration parameters.

\paragraph{Vote identifier namespace.}

The interval $[\configmax + 1, N - 1]$ is reserved for \emph{vote identifiers} (\voteid), which allows voters to verify that their vote has been included in the election state, without revealing or linking to the corresponding encrypted ballot. Vote identifiers are derived by hashing the process identifier, the voter’s address, and fresh randomness, and then mapping the result into the allowed range via modular reduction:
\[
\voteid = \texttt{voteIDMin} + \bigl(\text{Hash}(\processid,\address,k) \bmod (N - \texttt{voteIDMin})\bigr),
\]
where $\texttt{voteIDMin} = \texttt{configMax} + 1$. The \hlget{\statetransitioncircuit} recomputes this derivation and enforces that $\texttt{voteIDMin} \le \voteid < N$, guaranteeing that vote identifiers cannot overlap with configuration entries or ballot storage.

Unlike classical nullifiers, vote identifiers are not used to prevent duplicate voting directly. Instead, correctness is enforced at the circuit level: the state transition circuit guarantees that a vote identifier can only be inserted if the corresponding leaf in the state tree is empty. This emptiness check is enforced inside the ZK circuit. Since the vote identifier depends on a fresh random value $k$, if a collision occurs (i.e., the computed leaf is already occupied), the circuit will reject the transition. In that case, the voter must resample $k$ and recompute a new vote identifier until an unused leaf is obtained. As a result, collisions do not compromise correctness or liveness, but merely require retrying the identifier derivation.

Note that the maximum number of vote identifiers that can be stored is bounded by the size of this namespace ($\approx 3\cdot 2^{62}$), 
which upper-bounds the total number of votes and overwrites that can be processed throughout the election.
%This capacity is orders of magnitude larger than any realistic deployment, ensuring that identifier exhaustion is practically impossible.
This capacity is orders of magnitude larger than any realistic election workload: historical elections peak at $\approx 10^{9}$ votes~\cite{reuters2024india}, ensuring that identifier exhaustion is practically impossible.

\paragraph{Encrypted ballot namespace.}

The upper portion of the tree, $[N,\,2^{64}-1]$, is reserved for storing encrypted ballots. In contrast to vote identifiers, ballot storage locations are deterministic and derived from the voter's position in the census. Let $\idx$ denote the voter's census index, proven via a census membership proof. The storage path for the encrypted ballot is computed as
\[
\text{encPath} = N + 1 + \bigl(\idx \cdot 2^{16} + (\address \bmod 2^{16})\bigr).
\]
This construction assigns each eligible voter a unique, reserved leaf in the state tree. Because the index $\idx$ is authenticated by the census, a voter can only write to their own slot. This enables efficient \emph{last-vote-wins} semantics: if a voter overwrites their ballot, the new ciphertext replaces the previous one at the same path, while the tally is updated by subtracting the old contribution and adding the new one.

Together, the state tree provides a compact and tamper-proof commitment to the current status of the election. At the end of the election, the final state root together with the on-chain verification of all transitions serves as a complete cryptographic record of the election, encapsulating both the tally and the integrity of the entire voting procedure.